%%%%%%%%%%%%%%%%%%%%%%%%%%%%%%%%%%
%% Text formatting & manipulation
%%%%%%%%%%%%%%%%%%%%%%%%%%%%%%%%%%
\ExplSyntaxOn

%% Outputs a ``page \pageref{}'', optionally wrapped in braces
%% [t] = no wrapping, [c] = {curly}, [c] = (parenthesis), [b] = [braces]
\NewDocumentCommand{\RpgPage}{O{t} m}
{
	% O{t}: Optional argument #1 (default is 't' for standard 'text' style)
	% m: Mandatory argument #2 (the label)

	\def\temp@content{page~\pageref{#2}}

	\str_case:nnF {#1}
	{
		{p}{(\temp@content)}
			{b}{[\temp@content]}
			{c}{\{\temp@content\}}
			{t}{\temp@content}
	}
	{
		?\temp@content?
	}
}


%% Pluralises words based on the tokens
%% Arguments: [optional non-s based plural]{number}{word to pluralise}
%% Unstarred includes the number (\RpgPlural{3}{hat} = 3 hats)
%% Starred version just outputs the word (\RpgPlural*{1}{hat} = hat)
%% Complex plurals handles by the optional argument (\RpgPlural[octopodes]{5}{octopus} = 5 octopodes)
\NewDocumentCommand{\RpgPlural}{s o m m }
{
	%if unstarred, add no.
	\IfBooleanF{#1}
	{
		#3~
	}
	% #2: override plural (replaces +s) #3 count #4 word/phrase to be checked
	\int_compare:nTF {#3  = 1 }
	{
		#4
	}
	{
		\IfNoValueTF{#2}
		{
			#4s
		}
		{
			#2
		}
	}
}

% Take an integer and output the ordinal text
% \RpgOrdinal{1} = 1st, \RpgOrdinal{25} = 25th
% Optional first arg is a format applied to the ordinal text
\NewExpandableDocumentCommand{\RpgOrdinal}{o m}
{
	% Output the number first
	\int_to_arabic:n { #2 }
	%then the calculated suffix
	\IfNoValueTF{#1}
	{
		\__rpg_ComputeOrdinal:n { #2 }
	}
	{
		%Allow for arbitrary token insertion between the two (i.e. for formatting)
		#1{\__rpg_ComputeOrdinal:n { #2 }}
	}
}

% A direct wrapper for \RpgOrdinal[\textsuperscript]
\NewDocumentCommand{\RpgSuperOrdinal}{m}
{
	\RpgOrdinal[\textsuperscript]{#1}
}

%%%%%%%%%%%%%%%%%%%%%%%%%%%%%%%%%%
%% Internal functions for ordinal
%%%%%%%%%%%%%%%%%%%%%%%%%%%%%%%%%%

% Ordinal wrapper function, compute the suffix: checks for 11, 12, 13 first
\cs_new:Npn \__rpg_ComputeOrdinal:n #1
{

	% Check special cases 11, 12, 13 (modulo 100)
	\int_case:nnF {  \int_mod:nn { #1 } { 100 } }
	{
		{11} {th}
			{12} {th}
			{13} {th}
	}
	{
		% If not a teen, defer to the single-digit check
		\__rpg_ComputeNormalOrdinal:n { #1 }
	}
}

% compute the suffix based on the last digit
\cs_new:Npn \__rpg_ComputeNormalOrdinal:n #1
{
	% Get the last digit (modulo 10)

	\int_case:nnF{\int_mod:nn{#1}{10}}
	{
		{1} {st}
			{2} {nd}
			{3} {rd}
	}
	{
		% Default case
		{th}
	}
}


%%%%%%%%%%%%%%%%%%%
% Number formatter
%%%%%%%%%%%%%%%%%%%
%% Formats numbers with digit grouping (2300 -> 2,300), in the surrounding font. The options are given on each call,
%% rather than through \sisetup, so that a user's own siunitx commands are unaffected.
%% Commas in the input are ignored (so 2,300 is not read as 2.3), and anything which is not a number is printed as given
\tl_new:N \__rpg_numberSeparator
\tl_new:N \__rpg_numberGroupDigits

\NewDocumentCommand\RpgNumber{m}
{
	\siunitx_if_number:nTF {#1}
	{
		\num[
			input-decimal-markers = {.},
			input-ignore          = {,},
			group-separator       = {\__rpg_numberSeparator},
			group-minimum-digits  = {\__rpg_numberGroupDigits},
			mode                  = text,
			reset-text-family     = false,
			reset-text-series     = false,
			reset-text-shape      = false,
		]{#1}
	}
	{ #1 }
}
\NewDocumentCommand\RpgNumberSetSeparator{m}
{
	\tl_set:Nn \__rpg_numberSeparator {#1}
}
\NewDocumentCommand\RpgNumberSetDigits{m}
{
	\tl_set:Nn \__rpg_numberGroupDigits {#1}
}

\RpgNumberSetSeparator{,}
\RpgNumberSetDigits{4}
